\section{La crisis de los benchmarks: de puntuaciones comparables a señales de escenarios reales}

El segmento más sólido de esta conversación es el que trata sobre benchmark y evaluation. El punto de partida de Gu Yu (谷雨) es claro: lo que la academia y la industria necesitan de la evaluation nunca ha sido lo mismo. La academia quiere experimentos controlables para verificar hipótesis, definir problemas nuevos y comparar la capacidad de los métodos; la industria, en cambio, se preocupa más por si la mejora en la evaluación puede transferirse al despliegue real, si de verdad permite que el usuario complete la tarea, quede satisfecho, se quede o pague.

Esta distinción importa porque, en este momento, muchas personas dicen por un lado que la puntuación de benchmark no importa, y por otro lado no pueden prescindir del benchmark. Xie Tianbao (谢天宝) ofrece un juicio más matizado: el benchmark no es inútil, sino que depende de en qué etapa se encuentre. En la etapa de 0 a 1, puede definir el problema, exponer las carencias y demostrar que el modelo actual no lo logra; una vez que se entra en la etapa de optimización de 60 a 99, la puntuación pública cada vez es más susceptible a la síntesis de datos, al sobreajuste del conjunto de entrenamiento y a la personalización de la tarea. En ese momento, mirar solo la puntuación puede llevar a un juicio erróneo sobre la utilidad real.

\begin{importantbox}{El cambio de lugar del benchmark}
El benchmark funciona mejor como signal, no como final answer. Puede decirle a la comunidad «aquí hay un problema»«el modelo tiene una carencia aquí»«cierta capacidad está empezando a surgir»; pero no puede responder por sí solo «si este Agent es útil en un producto real». Cuanto más cerca se está de la implementación en la industria, más se necesitan usuarios reales, APIs reales, casos reales y un ciclo continuo de retroalimentación para calibrar.
\end{importantbox}

Shang Yanyi (尚晏仪) habla del reward a partir de los productos de viajes y así completa el eslabón más difícil para la industria. En un producto de viajes, al final se puede observar «si se vendió un boleto o no», pero ese reward está demasiado lejos de cada decisión que toma el modelo paso a paso. Que el usuario no compre el boleto puede deberse a que el precio no le convenció, a que no le gustó la ruta, a que no confía en la IA, o simplemente a que cambió de planes en el último momento. Que el usuario diga «está bien» tampoco equivale a que en verdad vaya a comprar. Es decir, el resultado final es muy claro, pero la atribución del proceso es muy difusa; hay mucha retroalimentación subjetiva, pero pocas señales entrenables.

Esto es distinto de la dificultad de los benchmark académicos. En las tareas académicas, el reward suele estar definido de antemano: se suman puntos si se acierta y se restan si se falla. Pero en un producto real, ni siquiera el usuario sabe con certeza qué experiencia es una buena experiencia. Por ejemplo, cuándo debe interrumpir un asistente de viajes por voz, cuándo debe escuchar, cuándo debe dar una explicación extensa: estos problemas no se pueden resolver con que un investigador escriba un benchmark guiado por su intuición. Que Shang Yanyi diga que dependen más de A/B Test y de una infraestructura de experimentación sostenible es, en realidad, reconocer que la preferencia del usuario no se escribe, sino que se revela poco a poco en la entrega real.

\begin{knowledgebox}{0 a 1 y 60 a 99}
\begin{tabularx}{\textwidth}{>{\raggedright\arraybackslash}p{0.2\textwidth} Y Y}
\toprule
Etapa & Valor del benchmark & Riesgo principal \\
\midrule
0 a 1 & Define tareas nuevas, demuestra que el modelo no puede hacerlas y construye un lenguaje común para el campo. & La tarea puede desligarse del escenario real y convertirse en «benchmark por el benchmark». \\
60 a 99 & Funciona como señal de una capacidad local y ayuda al equipo de ingeniería a ubicar sus carencias. & La puntuación se contamina fácilmente con el entrenamiento, el aumento de datos, la personalización de la tarea y la presentación selectiva. \\
Producto real & Debe usarse junto con la retroalimentación del usuario, los resultados comerciales y la atribución de fallas. & El reward es disperso, subjetivo y de cadena larga, y es difícil convertirlo directamente en señal de entrenamiento. \\
\bottomrule
\end{tabularx}
\end{knowledgebox}

El problema más amplio detrás de este segmento es que la investigación en Agent está pasando de «si la capacidad existe» a «cómo la capacidad produce valor de manera estable». En el primer caso, el benchmark es un instrumento científico; en el segundo, el benchmark es solo una parte del sistema de operación del producto. Lo verdaderamente difícil no es obtener una puntuación, sino conectar en un ciclo la puntuación, la retroalimentación del usuario, el análisis humano, el A/B Test, la selección del modelo, la estrategia de prompt, el flujo de herramientas y las métricas del negocio.

\begin{warningbox}{No atribuyas el fracaso de un producto a la puntuación de un solo modelo}
Cuando un producto de Agent fracasa, es muy probable que no sea algo tan simple como que «el modelo no es lo bastante fuerte». Tal vez la tarea se definió mal, tal vez el reward está demasiado lejos, tal vez la interacción no es confiable, o tal vez el usuario simplemente no quiere completar la tarea de esa manera. El benchmark solo puede iluminar una parte de las causas; no puede sustituir la atribución del producto.
\end{warningbox}

\subsection{Resumen de la sección}

La esencia de la crisis del benchmark no es que la puntuación pierda sentido, sino que se estrecha el margen de su interpretación. La academia todavía necesita el benchmark para definir problemas, mientras que la industria debe colocar la evaluation dentro del ciclo del producto real. La evaluación de los GUI Agent es especialmente difícil porque enfrenta tareas de cadena larga, experiencias subjetivas y reward disperso; la infrastructure que de verdad tiene valor es la infrastructure capaz de probar sin cesar, atribuir sin cesar y convertir sin cesar la retroalimentación del usuario en mejoras del sistema.

